\bookchapter{Kapitel 19: Fehlerbehebung und Selbsthilfe}

\chapterintro{Eine App stürzt beim Öffnen ab, oder sie friert plötzlich ein. Der Ablauf ist wie eine ärztliche Befragung – Schritt für Schritt vorgehen:}

\booksection{Zuerst schließen und neu öffnen}

: Die App aus der Übersicht der letzten Apps weg-wischen und erneut öffnen. Die meisten temporären Bugs sind damit behoben.

\booksubsection{App aktualisieren}

\begin{enumerate}
  \item \textbf{: Im App-Store nach Updates suchen. Manchmal liegt es daran, dass die App-Version zu alt ist und nicht zum System passt.}Berechtigungen prüfen
  \item \textbf{: Manche Apps benötigen eine bestimmte Berechtigung, um richtig zu funktionieren, aber du hast sie nicht gegeben. Zum Beispiel stürzt die Kamera-App ab, wenn ihr die Kamera-Berechtigung fehlt. Gehe in die Einstellungen, suche die App und gewähre alle nötigen Berechtigungen.}Cache leeren
  \item \textbf{: Einstellungen App-Verwaltung die App finden Speicher Cache leeren. Achtung: „Cache leeren", nicht „Daten löschen"! Daten zu löschen, entfernt den Login-Status und die Einstellungen der App – das entspricht einem Zurücksetzen auf Werkseinstellungen für die App.}Deinstallieren und neu installieren
  \item \textbf{: Der letzte Ausweg. Vor der Deinstallation sicherstellen, dass die App-Daten nicht verloren gehen – bei WeChat zuerst den Chatverlauf sichern, bei Spielen zuerst prüfen, ob der Account an die Handynummer gebunden ist.}III. Lädt nicht oder lädt langsam – vielleicht liegt es nicht am Akku
\end{enumerate}

\textbf{„Lädt nicht“ ist ein Problem, das besonders leicht Angst macht, aber oft liegt die Ursache bei den Zubehörteilen, nicht am defekten Smartphone.}Zuerst Ladeanschluss prüfen

\booksection{: Mit einer Zahnstange oder einem nicht leitenden kleinen Werkzeug vorsichtig Staub und Fussel aus dem Ladeanschluss herauskratzen. Das ist mir selbst passiert – ein Faden aus der Tasche hatte sich genau im Ladeanschluss verfangen, beim Einstecken dachte ich an Kontaktprobleme, nach dem Herausnehmen war alles gut.}

Anderes Kabel und anderes Ladegerät probieren

\begin{enumerate}
  \item \textbf{: Kabel haben eine viel höhere Ausfallrate als Ladegeräte, besonders oft geknickte Kabel. Ein definitiv intaktes Kabel und Ladegerät zum Kreuztest verwenden – schnell lässt sich feststellen, welches Zubehörteil defekt ist.}Problem mit kabellosem Laden
  \item \textbf{: Wenn du kabellos lädst, prüfe, ob die Smartphone-Hülle zu dick ist oder ob sich auf der Rückseite Metallteile befinden (magnetische Halterung o. Ä.). Außerdem: das Smartphone genau in die Mitte der Ladefläche legen – wenn es schief liegt, kann es sein, dass es nicht lädt.}Laden prüfen
  \item \textbf{: Bestätigen, dass die Leistung des Ladegeräts ausreicht (iPhone braucht für Schnellladen mindestens ein 20-W-PD-Ladegerät), dass das Kabel Schnellladen unterstützt (nicht alle Kabel können mehr als 3 A Strom führen), und dass das Smartphone nicht zu heiß ist (zu heiß automatische Leistungsreduktion zum Akkuschutz).}IV. Netzwerkprobleme – Wi-Fi will nicht
  \item \textbf{Draußen funktioniert mobile Daten einwandfrei, aber zu Hause lässt sich Wi-Fi nicht verbinden; oder Wi-Fi ist verbunden, aber es gibt kein Internet. In dieser Reihenfolge vorgehen:}Flugmodus ein/aus
  \item \textbf{: Kontrollzentrum nach unten ziehen, Flugmodus einschalten, 5 Sekunden warten, wieder ausschalten. Das ist, als würde das Smartphone neu nach dem Netz suchen – sehr effektiv.}Wi-Fi „vergessen" und neu verbinden
\end{enumerate}

\booksection{: In den Wi-Fi-Einstellungen auf das Ausrufezeichen neben dem Netzwerknamen tippen (bei Android das Zahnrad), „Dieses Netzwerk ignorieren“ / „Vergessen” wählen und dann das Passwort neu eingeben.}

Router neu starten

\begin{itemize}
  \item \textbf{: Viele Router zu Hause wurden monatelang nicht ausgeschaltet und werden mit der Zeit instabil. Strom ziehen, 30 Sekunden warten, wieder einstecken – die meisten Wi-Fi-Probleme sind gelöst.}Prüfen, ob das Smartphone das Problem ist
  \item \textbf{: Verbinden sich andere Geräte mit demselben Wi-Fi einwandfrei? Wenn ja, nur dein Smartphone nicht, muss es vielleicht neu gestartet werden oder die Netzwerkeinstellungen zurückgesetzt werden (Einstellungen Allgemein Übertragen oder zurücksetzen Netzwerkeinstellungen zurücksetzen). Das löscht keine Daten, aber alle Wi-Fi-Passwörter werden entfernt.}V. Bildschirmprobleme: Touch reagiert nicht, flackert, hinterlässt Abdrücke
  \item \textbf{Touch reagiert nicht}: Zuerst die Panzerglasfolie abziehen und sehen, ob sie schuld ist – billige Panzerglasfolien können nach längerer Nutzung tatsächlich den Touch beeinflussen. Wasser auf dem Bildschirm oder Öl an den Fingern kann ebenfalls zu Touch-Drift führen. Wenn beides nicht zutrifft, neu starten und erneut versuchen.
  \item \textbf{Abdrücke/Einbrennungen auf dem Bildschirm (Burn-in)}: Wenn lange Zeit dasselbe Bild angezeigt wird (z. B. ständig mit eingeschaltetem Display Douyin scrollen), kann das auf OLED-Bildschirmen „Einbrennungen" hinterlassen. Der Bildschirm ist nicht kaputt, das ist nur vorübergehend. Den Bildschirm eine halbe Stunde ruhen lassen, meist erholt er sich.
\end{itemize}

\booksection{Bildschirm flackert}

: Wenn bei geringer Helligkeit geflackert wird, kann es am PWM-Dimming liegen, häufig bei OLED-Bildschirmen. In den Einstellungen „DC-Dimming“ / „Anti-Flackern” suchen und aktivieren – danach wird es viel besser. Nicht alle Smartphones haben diese Funktion.

\begin{enumerate}
  \item \textbf{Ein Bereich des Bildschirms ist komplett schwarz/bleibt hell}: Das ist ernster – wahrscheinlich physische Beschädigung des Bildschirms, der ausgetauscht werden muss. Bei nicht-authorisierten Reparaturen sind Bildschirmfarbe und -helligkeit meist schlechter als beim Original, und die Wasserdichtigkeit leidet.
  \item \textbf{VI. Lässt sich nicht einschalten, hängt beim Logo – keine Panik}Diese Situation lässt einen am schnellsten das Herz schneller schlagen, aber ruhig bleiben – probiere zuerst Folgendes:
  \item \textbf{Eine halbe Stunde laden, dann versuchen}: Oft ist der Akku einfach komplett entleert. Original-Ladegerät anschließen, mindestens 30 Minuten warten, dann versuchen einzuschalten. Manche Smartphones reagieren nach vollständiger Entleerung erst, wenn eine bestimmte Ladung erreicht ist.
  \item \textbf{Erzwungener Neustart}: iPhone mit Full-Screen: einmal Lauter, einmal Leiser, dann die Seitentaste gedrückt halten, bis das Apple-Logo erscheint (evtl. 10–15 Sekunden gedrückt halten). Android-Methode: Ein-/Austaste länger als 10 Sekunden gedrückt halten, bis es vibriert oder der Bildschirm angeht.
\end{enumerate}

\booksection{Hängt beim Startbildschirm und startet ständig neu}

\begin{itemize}
  \item \textbf{: Das nennt sich „White Apple" (iPhone) oder „Logo hängen" (Android). Wenn der erzwungene Neustart nicht hilft, muss man evtl. in den Wiederherstellungsmodus und über den Computer reparieren. iPhone an den Computer anschließen, iTunes/Finder öffnen, in den Wiederherstellungsmodus wechseln und „Aktualisieren" wählen (nicht „Wiederherstellen") – das installiert das System neu, behält aber die Daten. Android in den Fastboot- oder Recovery-Modus wechseln; die Tastenkombination ist je nach Marke unterschiedlich – am besten „[Modell] Recovery-Modus" suchen.}Wenn nichts davon hilft
  \item \textbf{: Dann bleibt nur die professionelle Werkstatt. Wenn du ein Backup hast, musst du dir nicht zu viele Sorgen machen.}VII. Tonprobleme: Lautsprecher stumm, Anruflautstärke zu leise
  \item \textbf{Lautsprecher stumm}: Zuerst sicherstellen, dass es nicht der Stummschaltmodus ist (iPhone-Seitenschalter), nicht der Nicht-stören-Modus oder dass die Lautstärke ganz niedrig gestellt ist. Klingt dumm, aber das macht 40 \% der Fälle aus. Dann prüfen, ob Bluetooth mit einem Kopfhörer oder Lautsprecher verbunden ist.
  \item \textbf{Gesprächspartner hört dich nicht}: Prüfen, ob das Mikrofonloch verstopft ist – das kleine Loch unten am Smartphone verstopft leicht durch Staub. Mit einer weichen Bürste vorsichtig reinigen. Nicht mit Nadel oder Zahnstange hineinstechen – das verletzt das Staubgitter.
\end{itemize}

\booksection{Stimme kratzig/Verzerrung}

: Wenn der Lautsprecher Wasser bekommen hat – nicht mit dem Heißluftföhn trocknen! Mit Zellentuch die Oberfläche abtrocknen und dann an der Luft trocknen lassen. Im Netz gibt es „Wasser-ablassen-Apps“, die etwas bringen können (niedrige Frequenzen abspielen, um Wasser herauszuschütteln) – probier es aus.

\begin{itemize}
  \item \textbf{VIII. Autorisierte Werkstatt aufsuchen – diese Fallen vermeiden}Wenn alle oben genannten Methoden nicht geholfen haben, musst du in die professionelle Werkstatt. Aber der Reparaturmarkt ist tückisch – hier sind meine eigenen Erfahrungen und Lehren:
  \item \textbf{Die echte Autorisierung finden, nicht eine, die so aussieht}: Auf der Marken-Website (Apple.com.cn, Huawei.com.cn usw.) die Liste der „autorisierten Service-Center" suchen. Ein Laden am Straßenrand mit Apple-Logo, der nicht autorisiert ist, verwendet vielleicht keine Originalteile – nach der Reparatur ist deine Wasserdichtigkeit futsch.
  \item \textbf{Vor der Reparatur unbedingt tun}: Alle Daten sichern (verlass dich nicht darauf, dass bei der Reparatur nichts verloren geht), alle wichtigen Accounts ausloggen (Apple ID, WeChat, Alipay), „Mein Smartphone finden" deaktivieren. Manche Werkstätten verlangen das Sperrbildschirmpasswort – das ist ein Datenschutzrisiko. Am besten selbst vorher die Accounts ausloggen und dann abgeben.
  \item \textbf{Bildschirmwechsel}: Nicht-offizielle Bildschirme gibt es als „Original-gepresst" (Originalglas getauscht), „Nachgepresst" (nicht-original Glas) und „Zusammengebaut" (alles nicht-original). Die Preise unterscheiden sich um ein Vielfaches, aber auch Anzeigequalität und Haltbarkeit. Wenn das Smartphone noch in der Garantie ist, unbedingt den offiziellen Weg gehen.
\end{itemize}

\booksection{Akkuwechsel}

\begin{itemize}
  \item \textbf{: Nicht-original-Akkus haben oft übertriebene Kapazitätsangaben und verschlechtern sich nach einem halben Jahr stärker als ein Originalakku nach zwei Jahren. Bei Apple ist der offizielle Akkuwechsel nicht viel teurer als der inoffizielle – empfohlen, den offiziellen Weg zu gehen.}Falle der Datenwiederherstellung
  \item \textbf{: Wenn das Smartphone Wasser bekommen hat oder kaputtgefallen ist und die Werkstatt „Datenwiederherstellung kostet extra" sagt – Vorsicht. Manche unseriösen Betriebe übertreiben den Schaden, um mehr Geld zu verlangen. Seriöse Datenwiederherstellung bietet erst einen Kostenvoranschlag und nimmt bei Misserfolg kein Geld.}Der universelle erste Schritt: Smartphone neu starten
  \item \textbf{Egal welches Problem du hast – Ruckeln, Abstürze, kein Netz, Ton-Anomalien – starte zuerst das Smartphone neu. Das klingt wie eine Ausrede, aber nach meiner eigenen Erfahrung sind mindestens 60 \% der kleinen Probleme nach einem Neustart behoben. Das Smartphone ist wie ein Computer – wenn es lange läuft, sammeln sich im Speicher alle möglichen Dinge an, und ein Neustart leert alles und fängt neu an. Also: Wenn das Smartphone das nächste Mal ein Problem hat, nicht sofort „Was tun?" rufen – erst neu starten und dann schauen.}Präge dir jetzt ein: Diese Rettungsoperationen musst du auswendig kennen
\end{itemize}

\booksection{Erzwungener Neustart}

: Full-Screen-iPhone: Lauter Leiser Seitentaste gedrückt halten, bis das Logo erscheint. Android: Ein-/Austaste 10–15 Sekunden gedrückt halten, bis es vibriert.

\begin{itemize}
  \item \textbf{Ladeanschluss reinigen}: Ein kleines Werkzeug suchen (Zahnstange oder SIM-Nadel), vorsichtig den Staub aus dem Ladeanschluss entfernen.
  \item \textbf{Wi-Fi neu verbinden}: In den Wi-Fi-Einstellungen das aktuelle Wi-Fi „vergessen" und mit neu eingegebenem Passwort verbinden.
  \item \textbf{Flugmodus ein/aus}: Kontrollzentrum nach unten ziehen, Flugmodus einschalten, 5 Sekunden im Kopf zählen, ausschalten und prüfen, ob das Netz zurück ist.
  \item \textbf{Mache Screenshots dieser Operationsschritte oder schreibe sie in die Notizen-App. Wenn ein Problem auftritt, musst du nicht hektisch suchen – schau einfach in den Notizen nach.}Kapitel 20: Entwicklermodus und erweiterte Berechtigungen
  \item \textbf{Dieses Kapitel ist für Nutzer gedacht, die die „versteckten Funktionen" ihres Smartphones tiefer erkunden möchten. Entwicklermodus, Root-Rechte, BL-Sperre – diese sehr geekig klingenden Konzepte sind tatsächlich gar nicht so weit von normalen Nutzern entfernt.}I. Entwicklermodus: Versteckte Einstellungen freischalten
\end{itemize}

\begin{tcbwarn}{Der Entwicklermodus ist ein im Smartphone-System integriertes verstecktes Menü mit vielen Funktionen, die normale Nutzer nicht brauchen, die Entwickler jedoch benötigen.}
So aktivierst du ihn:

\end{tcbwarn}

\begin{tcbpractice}{Android}
\begin{enumerate}
  \item \textbf{Die echte Autorisierung finden, nicht eine, die so aussieht}HarmonyOS
  \item \textbf{: Einstellungen Über das Telefon 7-mal auf „HarmonyOS-Version" tippen.}iOS
  \item \textbf{: Normale Nutzer können den vollständigen Entwicklermodus nicht aktivieren – dafür ist ein kostenpflichtiges Apple-Developer-Konto nötig.}Häufige Funktionen:
  \item \textbf{USB-Debugging}: Verbindung zum Computer für fortgeschrittene Operationen wie Flashen, Datenwiederherstellung.
  \item Animationsgeschwindigkeit anpassen
\end{enumerate}

\end{tcbpractice}
